% Options for packages loaded elsewhere
\PassOptionsToPackage{unicode}{hyperref}
\PassOptionsToPackage{hyphens}{url}
\PassOptionsToPackage{dvipsnames,svgnames,x11names}{xcolor}
%
\documentclass[
  10.5pt,
]{article}
\usepackage{amsmath,amssymb}
\usepackage{iftex}
\ifPDFTeX
  \usepackage[T1]{fontenc}
  \usepackage[utf8]{inputenc}
  \usepackage{textcomp} % provide euro and other symbols
\else % if luatex or xetex
  \usepackage{unicode-math} % this also loads fontspec
  \defaultfontfeatures{Scale=MatchLowercase}
  \defaultfontfeatures[\rmfamily]{Ligatures=TeX,Scale=1}
\fi
\usepackage{lmodern}
\ifPDFTeX\else
  % xetex/luatex font selection
  \setmainfont[]{DejaVu Serif}
  \setmonofont[]{Noto Sans Mono}
\fi
% Use upquote if available, for straight quotes in verbatim environments
\IfFileExists{upquote.sty}{\usepackage{upquote}}{}
\IfFileExists{microtype.sty}{% use microtype if available
  \usepackage[]{microtype}
  \UseMicrotypeSet[protrusion]{basicmath} % disable protrusion for tt fonts
}{}
\makeatletter
\@ifundefined{KOMAClassName}{% if non-KOMA class
  \IfFileExists{parskip.sty}{%
    \usepackage{parskip}
  }{% else
    \setlength{\parindent}{0pt}
    \setlength{\parskip}{6pt plus 2pt minus 1pt}}
}{% if KOMA class
  \KOMAoptions{parskip=half}}
\makeatother
\usepackage{xcolor}
\usepackage[margin=0.8in]{geometry}
\usepackage{color}
\usepackage{fancyvrb}
\newcommand{\VerbBar}{|}
\newcommand{\VERB}{\Verb[commandchars=\\\{\}]}
\DefineVerbatimEnvironment{Highlighting}{Verbatim}{commandchars=\\\{\}}
% Add ',fontsize=\small' for more characters per line
\newenvironment{Shaded}{}{}
\newcommand{\AlertTok}[1]{\textcolor[rgb]{1.00,0.00,0.00}{\textbf{#1}}}
\newcommand{\AnnotationTok}[1]{\textcolor[rgb]{0.38,0.63,0.69}{\textbf{\textit{#1}}}}
\newcommand{\AttributeTok}[1]{\textcolor[rgb]{0.49,0.56,0.16}{#1}}
\newcommand{\BaseNTok}[1]{\textcolor[rgb]{0.25,0.63,0.44}{#1}}
\newcommand{\BuiltInTok}[1]{\textcolor[rgb]{0.00,0.50,0.00}{#1}}
\newcommand{\CharTok}[1]{\textcolor[rgb]{0.25,0.44,0.63}{#1}}
\newcommand{\CommentTok}[1]{\textcolor[rgb]{0.38,0.63,0.69}{\textit{#1}}}
\newcommand{\CommentVarTok}[1]{\textcolor[rgb]{0.38,0.63,0.69}{\textbf{\textit{#1}}}}
\newcommand{\ConstantTok}[1]{\textcolor[rgb]{0.53,0.00,0.00}{#1}}
\newcommand{\ControlFlowTok}[1]{\textcolor[rgb]{0.00,0.44,0.13}{\textbf{#1}}}
\newcommand{\DataTypeTok}[1]{\textcolor[rgb]{0.56,0.13,0.00}{#1}}
\newcommand{\DecValTok}[1]{\textcolor[rgb]{0.25,0.63,0.44}{#1}}
\newcommand{\DocumentationTok}[1]{\textcolor[rgb]{0.73,0.13,0.13}{\textit{#1}}}
\newcommand{\ErrorTok}[1]{\textcolor[rgb]{1.00,0.00,0.00}{\textbf{#1}}}
\newcommand{\ExtensionTok}[1]{#1}
\newcommand{\FloatTok}[1]{\textcolor[rgb]{0.25,0.63,0.44}{#1}}
\newcommand{\FunctionTok}[1]{\textcolor[rgb]{0.02,0.16,0.49}{#1}}
\newcommand{\ImportTok}[1]{\textcolor[rgb]{0.00,0.50,0.00}{\textbf{#1}}}
\newcommand{\InformationTok}[1]{\textcolor[rgb]{0.38,0.63,0.69}{\textbf{\textit{#1}}}}
\newcommand{\KeywordTok}[1]{\textcolor[rgb]{0.00,0.44,0.13}{\textbf{#1}}}
\newcommand{\NormalTok}[1]{#1}
\newcommand{\OperatorTok}[1]{\textcolor[rgb]{0.40,0.40,0.40}{#1}}
\newcommand{\OtherTok}[1]{\textcolor[rgb]{0.00,0.44,0.13}{#1}}
\newcommand{\PreprocessorTok}[1]{\textcolor[rgb]{0.74,0.48,0.00}{#1}}
\newcommand{\RegionMarkerTok}[1]{#1}
\newcommand{\SpecialCharTok}[1]{\textcolor[rgb]{0.25,0.44,0.63}{#1}}
\newcommand{\SpecialStringTok}[1]{\textcolor[rgb]{0.73,0.40,0.53}{#1}}
\newcommand{\StringTok}[1]{\textcolor[rgb]{0.25,0.44,0.63}{#1}}
\newcommand{\VariableTok}[1]{\textcolor[rgb]{0.10,0.09,0.49}{#1}}
\newcommand{\VerbatimStringTok}[1]{\textcolor[rgb]{0.25,0.44,0.63}{#1}}
\newcommand{\WarningTok}[1]{\textcolor[rgb]{0.38,0.63,0.69}{\textbf{\textit{#1}}}}
\usepackage{longtable,booktabs,array}
\usepackage{calc} % for calculating minipage widths
% Correct order of tables after \paragraph or \subparagraph
\usepackage{etoolbox}
\makeatletter
\patchcmd\longtable{\par}{\if@noskipsec\mbox{}\fi\par}{}{}
\makeatother
% Allow footnotes in longtable head/foot
\IfFileExists{footnotehyper.sty}{\usepackage{footnotehyper}}{\usepackage{footnote}}
\makesavenoteenv{longtable}
\setlength{\emergencystretch}{3em} % prevent overfull lines
\providecommand{\tightlist}{%
  \setlength{\itemsep}{0pt}\setlength{\parskip}{0pt}}
\setcounter{secnumdepth}{5}
\ifLuaTeX
  \usepackage{selnolig}  % disable illegal ligatures
\fi
\usepackage{bookmark}
\IfFileExists{xurl.sty}{\usepackage{xurl}}{} % add URL line breaks if available
\urlstyle{same}
\hypersetup{
  pdftitle={Pseudo-Last Design for SUPR-Guided Shoe Reconstruction},
  pdfauthor={Prepared for Abhinay Belde},
  colorlinks=true,
  linkcolor={blue},
  filecolor={Maroon},
  citecolor={Blue},
  urlcolor={blue},
  pdfcreator={LaTeX via pandoc}}

\title{Pseudo-Last Design for SUPR-Guided Shoe Reconstruction}
\author{Prepared for Abhinay Belde}
\date{May 31, 2026}

\begin{document}
\maketitle

{
\hypersetup{linkcolor=}
\setcounter{tocdepth}{3}
\tableofcontents
}
\section{Pseudo-Last Design for SUPR-Guided Shoe
Reconstruction}\label{pseudo-last-design-for-supr-guided-shoe-reconstruction}

\subsection{Design target}\label{design-target}

For your use case, the right target is \textbf{not} an industrial
manufacturing last and \textbf{not} a literal copy of the SUPR foot. It
is a \textbf{computational, watertight, smooth interior prior} that
preserves the hidden sole/interior support region well enough to
regularize a NeuS-style or G-Shell-style reconstruction when the shoe
opening hides much of the inner topology. That framing matches the
literature on both sides of your problem: NeuS reconstructs a zero-level
SDF from images but notes that severe self-occlusion and thin structures
remain difficult without stronger geometric constraints, while G-Shell
explicitly represents open surfaces as islands on a watertight shell
using a manifold SDF, making a watertight interior-support prior
especially natural for the hidden parts of a shoe. SUPR is also a better
anatomical seed than older generic body models because it explicitly
improves feet, toe motion, and foot deformation under ground contact
using large-scale body-part and 4D foot scans.

In digital footwear work, the shoe last is treated as the core
determinant of shoe shape, fit, and volume, and custom-last workflows
typically align the foot to a heel-based axis, then apply
\textbf{region-specific ease allowances} and \textbf{cross-sectional
editing} rather than copying raw anatomy vertex-for-vertex. That is
exactly the right design philosophy here: treat the SUPR foot as the
anatomical low-frequency prior, and build a smoother, more shoe-like
pseudo-last that is intentionally biased toward stable hidden topology.

My recommendation is therefore:

\textbf{Use a section-based pseudo-last pipeline.}\\
Preserve the plantar support geometry and overall hindfoot/forefoot
volume cues from the aligned SUPR foot, but replace the literal
dorsal/toe anatomy with a smoother last-like upper envelope, then
convert that mesh into a watertight SDF/occupancy prior for shell and
mSDF supervision.

\subsection{Why a last is not a foot}\label{why-a-last-is-not-a-foot}

A real shoe last is a \textbf{rigid shaping tool}, not a literal foot
twin. Footwear references and last-design papers describe the last as
the object that determines shoe style, fit, and internal volume, and
last-processing papers explicitly note that a scanned foot must
typically be \textbf{adjusted} before it matches last conventions
because footwear includes features such as toe spring, heel height, and
intentionally distributed ease. In other words, the last is already a
geometric abstraction of the foot, with added design bias.

A foot, by contrast, is articulated, compliant, and dynamically
changing. Dynamic foot-shape studies show that foot morphology changes
under loading; safety-shoe and walking studies report practically
relevant changes in width and girth between static and dynamic
conditions, and 4D statistical models capture both inter-person and
within-gait variation. That matters because copying the bare anatomical
foot too faithfully will overfit the wrong signal in exactly the regions
that should instead become a smoother shoe cavity.

The forefoot makes this especially clear. Toe-box studies show that shoe
toe-box shape and volume materially affect dorsal and plantar forefoot
pressures, with sharper or more constrained toe regions increasing
pressure in sensitive areas. Put bluntly: \textbf{individual toes are
anatomically real, but they are not the right topology for a sneaker
toe-box prior}.

A final point is population variation. Large-scale foot-scan analyses
and footwear-focused reviews show substantial differences across sex and
region, and the review literature is explicit that women's feet should
not be treated as scaled-down men's feet. For your problem, the
implication is not that you need sex-specific last engineering on day
one, but that a good pseudo-last should be \textbf{parametric and
smoothable}, not hard-coded around one raw foot surface.

\subsection{Anatomy to preserve and anatomy to
modify}\label{anatomy-to-preserve-and-anatomy-to-modify}

The simplest useful rule is this:

\textbf{Preserve low-frequency, support-relevant anatomy. Smooth
high-frequency, shoe-hostile anatomy.}

The parts I would preserve are the \textbf{plantar heel}, \textbf{ball
region}, \textbf{general forefoot spread}, \textbf{overall instep dome},
\textbf{heel bulb / calcaneal back volume}, and a \textbf{moderated
medial arch uplift}. Those are the features most closely tied to
support, contact, and functional fit in the footwear and orthotic
literature, and they are exactly the kinds of landmarks and measurements
most foot-shape pipelines focus on: heel point, toe point, ball
landmarks, arch region, dorsal height, midfoot width, and girths.
Plantar-shape studies also show that weight-bearing plantar geometry
matters for support design, so preserving bottom shape in some form is
important.

The parts I would deliberately smooth or remove are the
\textbf{shank/leg above the ankle}, \textbf{separate toes},
\textbf{sharp metatarsal-head bumps}, \textbf{malleolar/ankle
prominences}, and any small dorsum/navicular protrusions that create
local waviness without adding useful support semantics for a closed
sneaker. That choice is consistent with section-based last workflows and
with ease-allowance approaches that distribute fit volume smoothly
across last regions rather than preserving every anatomical protrusion.

So, in practical terms:

\begin{itemize}
\tightlist
\item
  \textbf{Keep:} plantar support shape, hindfoot cup volume, forefoot
  width trend, instep height trend.
\item
  \textbf{Soften:} arch details into a clean midfoot sweep.
\item
  \textbf{Merge:} toes into one toe-box mass.
\item
  \textbf{Tighten:} upper heel slightly for hold.
\item
  \textbf{Delete:} anything above the intended collar/opening zone.
\end{itemize}

That gives you a prior that is anatomically anchored where it matters
and topologically conservative where image supervision is weak.

\subsection{Mesh construction
algorithm}\label{mesh-construction-algorithm}

\subsubsection{Preprocess and crop the SUPR
foot}\label{preprocess-and-crop-the-supr-foot}

Use the aligned neutral SUPR right foot as the starting mesh \(M_f\).
Because your shoes are already normalized, operate entirely in
normalized coordinates and define the foot length

\[
L_f = x_{\max}(M_f) - x_{\min}(M_f).
\]

I recommend estimating a \textbf{leg axis} from the rear-top vertex
cluster rather than using a naive axis-aligned clip. If SUPR joint data
are available, use the ankle joint and local shank axis directly;
otherwise, estimate the axis by PCA over vertices in the rearfoot and
upper-shank cluster. Then sweep a family of planes orthogonal to that
axis near the ankle and pick the plane whose section perimeter is
minimal. This removes the leg above the ankle much more cleanly than
clipping by \(x=\text{const}\). That is consistent with footwear
alignment literature, which emphasizes robust reference axes and warns
that simple heel-to-toe alignment can be unstable when toe geometry is
irregular.

After that crop, keep the foot-side component and triangulate a
temporary cap over the cut boundary. This temporary cap is only for
downstream section extraction and SDF cleanup; it does not need to be
final-quality.

\subsubsection{Build a shoe-constrained centerline and width
profile}\label{build-a-shoe-constrained-centerline-and-width-profile}

Even if you already have a support centerline, I would recompute a
\textbf{sectional centerline} from the support footprint as a
consistency check.

At each slice \(x_i\), intersect the support footprint with the line
\(x=x_i\). Let the largest connected interval be

\[
I_i = [z_L(x_i), z_R(x_i)].
\]

Define the raw center and support width as

\[
c_i^0 = \frac{z_L(x_i) + z_R(x_i)}{2}, \qquad
w_i^0 = z_R(x_i)-z_L(x_i).
\]

Then smooth both with cubic smoothing splines or a discrete
second-derivative regularizer:

\[
c = \arg\min_{\tilde c}\sum_i \left(\tilde c(x_i)-c_i^0\right)^2 + \lambda_c \sum_i (\Delta^2 \tilde c_i)^2,
\]

\[
w = \arg\min_{\tilde w}\sum_i \left(\tilde w(x_i)-w_i^0\right)^2 + \lambda_w \sum_i (\Delta^2 \tilde w_i)^2.
\]

Because your shoes are already in a heel-to-toe coordinate frame, this
x-monotone formulation is simpler and more stable than a full
medial-axis skeleton for normal sneakers. Foot-measurement and
retail-scan pipelines use essentially this kind of axis/outline logic
when extracting robust length, width, and heel-width measurements.

\subsubsection{Preserve the plantar shape using the footbed and the SUPR
arch}\label{preserve-the-plantar-shape-using-the-footbed-and-the-supr-arch}

Let \(B(x,z)\) be your pseudo-footbed surface in the given coordinates,
where larger \(y\) means more downward/closer to the sole. Let
\(y_f^{pl}(x,z)\) be the plantar surface of the cropped SUPR foot after
local interpolation.

You do \textbf{not} want to flatten the pseudo-last bottom to the
footbed everywhere, because that erases useful arch information. But you
also do \textbf{not} want to use the bare foot plantar surface directly,
because that may float unrealistically relative to the shoe support
geometry.

A useful blend is:

\[
y_b(x,z) = B(x,z) - \rho_b(s)\,[B(x,z)-y_f^{pl}(x,z)]_+,
\qquad s=\frac{x-x_h}{L_f},
\]

where \([u]_+=\max(u,0)\), \(x_h=x_{\min}(M_f)\), and \(\rho_b(s)\) is
an arch-preservation weight that is strong in the midfoot and weak in
heel/toe support zones.

A good default is a Gaussian-like arch weight:

\[
\rho_b(s)=\rho_{\text{arch}}\exp\!\left(-\frac{(s-0.45)^2}{2(0.12)^2}\right).
\]

This gives you a bottom surface that follows the shoe-derived footbed in
supported regions while retaining moderated arch lift from the
anatomical foot in the midfoot. Plantar-shape studies under different
weight-bearing conditions strongly support the idea that the plantar
surface should be treated as a controlled design variable in anything
support-related.

\subsubsection{Extract cross-sections and merge the
toes}\label{extract-cross-sections-and-merge-the-toes}

Now slice the cropped foot mesh at \(N_x\) equally spaced \(x\)-planes.
Equal-interval sectioning is standard in digital last design; classic
foot-to-last work even describes 1\% foot-length sections, and more
recent CAD work still builds shoe lasts from regular cross-sections and
characteristic points.

For each slice \(x_i\):

\begin{enumerate}
\def\labelenumi{\arabic{enumi}.}
\tightlist
\item
  Intersect the mesh with the plane \(x=x_i\).
\item
  Convert the resulting contour to local slice coordinates \((\xi, h)\),
  where \[
  \xi = z-c(x_i), \qquad h = y_b(x_i,z)-y.
  \] Here \(h\) is height above the last bottom.
\item
  Rasterize the slice to a binary mask \(M_i(\xi,h)\).
\end{enumerate}

For the toe region, do \textbf{not} use the raw mask directly. Instead,
starting at a normalized longitudinal position \(s_{\text{merge}}\),
apply morphological closing with an elliptical kernel whose radii
increase toward the toe tip:

\[
M_i' = (M_i \oplus K_i)\ominus K_i.
\]

This closes the inter-toe gaps and turns the literal toe anatomy into
one continuous toe-box section. It is the simplest robust way to enforce
``one toe-box volume'' without hand-sculpting. After closing, extract
the contour of \(M_i'\) and fit a smooth superellipse or resampled
B-spline.

\subsubsection{Add toe allowance, heel hold, and optional side/top
inflation}\label{add-toe-allowance-heel-hold-and-optional-sidetop-inflation}

Compute a smoothed raw foot half-width \(a_f(x)\) and raw foot height
\(H_f(x)\) from the post-processed slice masks, then regularize them
along \(x\) with the same second-derivative smoothing used for the
centerline.

Then define target allowances:

\[
\delta_x = \alpha_t L_f
\]

for toe length,

\[
\delta_w(s)=W_f\left(\alpha_{mid} g(s;0.30,0.55) + \alpha_{fore} g(s;0.60,0.85)\right)
\]

for side clearance, and

\[
\delta_h(s)=H_f(s)\left(\beta_{vamp} g(s;0.35,0.65)+\beta_{toe} g(s;0.72,0.95)\right)
\]

for top clearance, where \(g(s;a,b)\) is a smoothstep profile.

Then clamp the width to the shoe support geometry:

\[
a_t(x)= \min\!\left(\eta_s\,\frac{w(x)}{2},\; a_f(x)+\delta_w(s)\right),
\]

\[
H(x)=H_f(x)+\delta_h(s).
\]

The factor \(\eta_s<1\) is important when your support profile comes
from an outer shoe mesh rather than the true inner cavity.

For the upper envelope, use a superellipse section:

\[
H_{\text{shoe}}(x,\xi)=H(x)\left[1-\left|\frac{\xi}{a_t(x)}\right|^{p(s)}\right]_+^{1/q(s)},
\]

\[
y_u(x,z) = y_b(x,z) - H_{\text{shoe}}(x, z-c(x)).
\]

Then add a \textbf{height-dependent heel-hold contraction} in the
hindfoot by shrinking the slice laterally after section generation:

\[
z \leftarrow c(x) + \big(z-c(x)\big)\Big(1-\gamma_h(s)\tau^\nu\Big),
\qquad
\tau=\frac{y_b(x,z)-y}{H(x)},
\]

with \(\gamma_h(s)\) nonzero only in the hindfoot. This keeps the
plantar heel width stable while tightening the upper heel volume, which
is much closer to what a sneaker backpart actually needs.

For toe allowance, append a short extension cap of length \(\delta_x\)
beyond the anatomical toe tip and taper width/height smoothly to zero,
for example with a hemispheroidal taper:

\[
a_t(x_t+u)=a_t(x_t)\sqrt{1-(u/\delta_x)^2},
\qquad
H(x_t+u)=H(x_t)\sqrt{1-(u/\delta_x)^2},
\quad u\in[0,\delta_x].
\]

\subsubsection{Loft, cap, remesh, and do an SDF cleanup
round-trip}\label{loft-cap-remesh-and-do-an-sdf-cleanup-round-trip}

After generating all contours:

\begin{itemize}
\tightlist
\item
  resample every contour to the same point count \(N_\theta\),
\item
  loft adjacent slices with a strip of triangles,
\item
  triangulate the toe cap,
\item
  triangulate the ankle-opening cap from the crop contour,
\item
  isotropically remesh,
\item
  smooth only non-plantar vertices,
\item
  then do one \textbf{SDF round-trip}: mesh \(\rightarrow\) signed grid
  \(\rightarrow\) Marching Cubes.
\end{itemize}

That last step is the easiest way to remove small self-intersections and
guarantee a watertight output. Because this pseudo-last is
computational, not manufacturing-grade, the SDF round-trip is a feature,
not a bug.

\subsubsection{Reference pseudocode}\label{reference-pseudocode}

\begin{Shaded}
\begin{Highlighting}[]
\KeywordTok{def}\NormalTok{ build\_pseudolast(}
\NormalTok{    supr\_mesh,}
\NormalTok{    foot\_transform,}
\NormalTok{    footbed\_B,           }\CommentTok{\# callable B(x, z)}
\NormalTok{    support\_footprint,   }\CommentTok{\# x{-}z mask or interval function}
\NormalTok{    N\_x}\OperatorTok{=}\DecValTok{128}\NormalTok{,}
\NormalTok{    N\_theta}\OperatorTok{=}\DecValTok{128}
\NormalTok{):}
\NormalTok{    M }\OperatorTok{=}\NormalTok{ transform\_mesh(supr\_mesh, foot\_transform)}

    \CommentTok{\# {-}{-}{-} crop leg above ankle {-}{-}{-}}
\NormalTok{    leg\_axis }\OperatorTok{=}\NormalTok{ estimate\_leg\_axis\_pca(M)}
\NormalTok{    crop\_plane }\OperatorTok{=}\NormalTok{ choose\_min\_perimeter\_ankle\_plane(M, leg\_axis)}
\NormalTok{    M }\OperatorTok{=}\NormalTok{ cut\_keep\_foot\_side(M, crop\_plane)}

    \CommentTok{\# {-}{-}{-} support centerline / width {-}{-}{-}}
\NormalTok{    xs }\OperatorTok{=}\NormalTok{ np.linspace(xmin(M), xmax(M), N\_x)}
\NormalTok{    c\_raw, w\_raw }\OperatorTok{=}\NormalTok{ sectional\_centerline\_width(support\_footprint, xs)}
\NormalTok{    c }\OperatorTok{=}\NormalTok{ smooth\_profile(c\_raw, lam}\OperatorTok{=}\FloatTok{1e{-}3}\NormalTok{)}
\NormalTok{    w }\OperatorTok{=}\NormalTok{ smooth\_profile(w\_raw, lam}\OperatorTok{=}\FloatTok{1e{-}3}\NormalTok{)}

    \CommentTok{\# {-}{-}{-} plantar bottom from footbed + moderated arch preservation {-}{-}{-}}
\NormalTok{    y\_bottom }\OperatorTok{=}\NormalTok{ build\_bottom\_surface(M, footbed\_B, c, arch\_strength}\OperatorTok{=}\FloatTok{0.75}\NormalTok{)}

    \CommentTok{\# {-}{-}{-} extract and clean sections {-}{-}{-}}
\NormalTok{    sections }\OperatorTok{=}\NormalTok{ []}
    \ControlFlowTok{for}\NormalTok{ x }\KeywordTok{in}\NormalTok{ xs:}
\NormalTok{        contour }\OperatorTok{=}\NormalTok{ mesh\_section(M, x)}
\NormalTok{        mask }\OperatorTok{=}\NormalTok{ rasterize\_section(contour, center\_z}\OperatorTok{=}\NormalTok{c(x), bottom}\OperatorTok{=}\NormalTok{y\_bottom)}
\NormalTok{        s }\OperatorTok{=}\NormalTok{ (x }\OperatorTok{{-}}\NormalTok{ xs[}\DecValTok{0}\NormalTok{]) }\OperatorTok{/}\NormalTok{ (xs[}\OperatorTok{{-}}\DecValTok{1}\NormalTok{] }\OperatorTok{{-}}\NormalTok{ xs[}\DecValTok{0}\NormalTok{])}

        \ControlFlowTok{if}\NormalTok{ s }\OperatorTok{\textgreater{}=} \FloatTok{0.74}\NormalTok{:}
\NormalTok{            mask }\OperatorTok{=}\NormalTok{ morphological\_close(mask, kernel}\OperatorTok{=}\NormalTok{toe\_kernel(s))}

\NormalTok{        a\_f, H\_f }\OperatorTok{=}\NormalTok{ fit\_section\_stats(mask)}
\NormalTok{        a\_t, H\_t }\OperatorTok{=}\NormalTok{ apply\_clearances(}
\NormalTok{            a\_f}\OperatorTok{=}\NormalTok{a\_f,}
\NormalTok{            H\_f}\OperatorTok{=}\NormalTok{H\_f,}
\NormalTok{            support\_half\_width}\OperatorTok{=}\FloatTok{0.93} \OperatorTok{*} \FloatTok{0.5} \OperatorTok{*}\NormalTok{ w(x),}
\NormalTok{            s}\OperatorTok{=}\NormalTok{s}
\NormalTok{        )}

\NormalTok{        contour }\OperatorTok{=}\NormalTok{ superellipse\_contour(a\_t, H\_t, p}\OperatorTok{=}\NormalTok{section\_p(s), q}\OperatorTok{=}\NormalTok{section\_q(s))}
\NormalTok{        contour }\OperatorTok{=}\NormalTok{ apply\_heel\_hold(contour, s}\OperatorTok{=}\NormalTok{s, shrink}\OperatorTok{=}\FloatTok{0.05}\NormalTok{)}
\NormalTok{        contour }\OperatorTok{=}\NormalTok{ map\_section\_to\_world(contour, x}\OperatorTok{=}\NormalTok{x, center\_z}\OperatorTok{=}\NormalTok{c(x), bottom}\OperatorTok{=}\NormalTok{y\_bottom)}
\NormalTok{        sections.append(contour)}

\NormalTok{    sections }\OperatorTok{=}\NormalTok{ add\_toe\_extension\_cap(sections, toe\_allowance}\OperatorTok{=}\FloatTok{0.05} \OperatorTok{*}\NormalTok{ foot\_length(M))}
\NormalTok{    mesh }\OperatorTok{=}\NormalTok{ loft\_sections(sections)}
\NormalTok{    mesh }\OperatorTok{=}\NormalTok{ cap\_openings(mesh)         }\CommentTok{\# ankle cap + remaining caps}
\NormalTok{    mesh }\OperatorTok{=}\NormalTok{ smooth\_nonplantar(mesh)}
\NormalTok{    mesh }\OperatorTok{=}\NormalTok{ sdf\_roundtrip(mesh, resolution}\OperatorTok{=}\DecValTok{192}\NormalTok{)}

    \ControlFlowTok{return}\NormalTok{ mesh}
\end{Highlighting}
\end{Shaded}

\subsection{Mathematical formulation and normalized
defaults}\label{mathematical-formulation-and-normalized-defaults}

\subsubsection{Coordinates and section
model}\label{coordinates-and-section-model}

Use your coordinate convention directly:

\begin{itemize}
\tightlist
\item
  \(+X\): heel \(\rightarrow\) toe
\item
  \(+Y\): sole/down
\item
  \(-Y\): opening/up
\item
  \(+Z\): width
\end{itemize}

Define the normalized longitudinal coordinate

\[
s = \frac{x-x_h}{L_f}, \qquad L_f=x_t-x_h.
\]

Define height above the bottom reference using your footbed:

\[
h(x,y,z) = y_b(x,z)-y,
\]

where \(y_b\) is the blended plantar surface from the previous section.

With this convention, the pseudo-last volume can be represented
slice-wise as

\[
\Omega_{\text{last}}=\Big\{(x,y,z)\;:\; x\in[x_h,\,x_t+\delta_x],\;
y_u(x,z)\le y \le y_b(x,z)\Big\},
\]

and \(y_u\) is the smoother shoe-like upper envelope built from section
width and height profiles.

The profiles that matter most are:

\begin{itemize}
\tightlist
\item
  centerline \(c(x)\),
\item
  half-width \(a_t(x)\),
\item
  section height \(H(x)\),
\item
  arch-preservation weight \(\rho_b(s)\),
\item
  heel-hold strength \(\gamma_h(s)\).
\end{itemize}

All of them should be regularized along \(x\) by curvature penalties,
not just smoothed ad hoc. In practice, the discrete objective

\[
f^{sm}=\arg\min_f \sum_i (f_i-\hat f_i)^2 + \lambda \sum_i (\Delta^2 f_i)^2
\]

works well for \(f\in\{c,w,a_f,H_f\}\).

\subsubsection{Recommended normalized
priors}\label{recommended-normalized-priors}

Because your meshes are normalized, I would treat all defaults as
\textbf{dimensionless fractions} of foot length, foot width, or local
sectional height. The toe-length numbers below are anchored to
footwear-fit literature that commonly recommends about 10-15 mm of extra
length for adults; the width-related values are engineering priors
informed by the literature on dynamic foot widening, pliable fit, and
the fact that footwear fit still lacks one universal quantitative
standard.

\begin{longtable}[]{@{}
  >{\raggedright\arraybackslash}p{(\columnwidth - 6\tabcolsep) * \real{0.2000}}
  >{\raggedleft\arraybackslash}p{(\columnwidth - 6\tabcolsep) * \real{0.2667}}
  >{\raggedleft\arraybackslash}p{(\columnwidth - 6\tabcolsep) * \real{0.2667}}
  >{\raggedleft\arraybackslash}p{(\columnwidth - 6\tabcolsep) * \real{0.2667}}@{}}
\toprule\noalign{}
\begin{minipage}[b]{\linewidth}\raggedright
Quantity
\end{minipage} & \begin{minipage}[b]{\linewidth}\raggedleft
Symbol
\end{minipage} & \begin{minipage}[b]{\linewidth}\raggedleft
Recommended range
\end{minipage} & \begin{minipage}[b]{\linewidth}\raggedleft
Good default
\end{minipage} \\
\midrule\noalign{}
\endhead
\bottomrule\noalign{}
\endlastfoot
Toe length allowance & \(\alpha_t\) & \(0.040L_f\) to \(0.060L_f\) &
\(0.050L_f\) \\
Toe merge start & \(s_{\text{merge}}\) & 0.72 to 0.78 & 0.74 \\
Full toe-box by & \(s_{\text{box}}\) & 0.86 to 0.92 & 0.88 \\
Midfoot side clearance & \(\alpha_{mid}\) & \(0.010W_f\) to \(0.030W_f\)
& \(0.020W_f\) \\
Forefoot side clearance & \(\alpha_{fore}\) & \(0.030W_f\) to
\(0.060W_f\) & \(0.040W_f\) \\
Vamp / top clearance & \(\beta_{vamp}\) & 0.04 to 0.08 of local height &
0.05 \\
Toe vertical clearance & \(\beta_{toe}\) & 0.03 to 0.07 of local height
& 0.04 \\
Heel clearance & \(\alpha_{heel}\) & 0.01 to 0.03 of local heel width &
0.015 \\
Heel-hold upper shrink & \(\gamma_h\) & 0.03 to 0.08 of hindfoot upper
width & 0.05 \\
Arch/bottom preservation & \(\rho_{\text{arch}}\) & 0.60 to 0.90 &
0.75 \\
Support-to-inner width scale & \(\eta_s\) & 0.90 to 0.96 & 0.93 \\
Forefoot superellipse exponent & \(p\) & 3.0 to 4.5 & 3.5 \\
Vertical superellipse exponent & \(q\) & 1.8 to 2.8 & 2.2 \\
\end{longtable}

A few practical notes matter here.

First, \textbf{toe allowance should be longer than side allowance is
wide}. If you inflate width aggressively but barely extend length, the
pseudo-last becomes bulbous rather than sneaker-like.

Second, \textbf{heel clearance should stay meaningfully smaller than
forefoot side clearance}. Otherwise the prior will teach the shell that
heel looseness is acceptable, which is usually the opposite of what
stabilizes topology in a low-top shoe.

Third, if your support width was extracted from the \textbf{outer} shoe
mesh, start with \(\eta_s \approx 0.93\). If it came from a cleaner
estimate of the footbed/interior support region, you can move it toward
0.96.

\subsection{SDF conversion and mSDF-guided
training}\label{sdf-conversion-and-msdf-guided-training}

\subsubsection{Convert the pseudo-last mesh to an SDF or occupancy
field}\label{convert-the-pseudo-last-mesh-to-an-sdf-or-occupancy-field}

Once you have a watertight pseudo-last mesh \(M_L\), convert it to a
continuous field on the same domain as your reconstruction model.

For any query point \(p\), define the unsigned distance

\[
d_u(p) = \min_{t\in M_L} \|p-\Pi_t(p)\|_2,
\]

where \(\Pi_t(p)\) is the closest point on triangle \(t\). For the sign,
if the mesh is watertight, ray parity is fine; if you want robustness to
small imperfections after lofting, use a generalized or fast
winding-number test. That is a standard geometry-processing choice and
works well for voxelization and signing distances.

Then define

\[
d_L(p)=\text{sign}(p)\,d_u(p),
\]

and a truncated field

\[
\tilde d_L(p)=\mathrm{clip}(d_L(p),-\tau,\tau),
\qquad \tau \in [0.03L_f,\;0.06L_f].
\]

If you need occupancy instead of SDF, use either

\[
o_L(p)=\mathbf{1}[d_L(p)\le 0]
\]

or a differentiable soft occupancy

\[
\hat o_L(p)=\sigma\!\left(-\frac{d_L(p)}{\sigma_o}\right),
\qquad \sigma_o \approx 0.01L_f.
\]

DeepSDF and Occupancy Networks are still the right conceptual references
here: continuous SDFs give you a boundary plus inside/outside semantics,
while occupancy fields give you a smooth classifier at arbitrary
resolution.

\subsubsection{Use the pseudo-last in shell training
first}\label{use-the-pseudo-last-in-shell-training-first}

For G-Shell-style training, I would use the pseudo-last in \textbf{two
stages}.

The first stage regularizes the \textbf{watertight shell SDF}
\(d_\theta\). Assume the usual sign convention that \(d_\theta(p)<0\)
means the point lies inside the current shell volume.

Sample points on the pseudo-last surface \(S_L\), and enforce that the
pseudo-last sits strictly inside the learned shell by a margin
\(m_{in}\):

\[
\mathcal L_{\text{contain}}
=
\mathbb E_{p\sim S_L}
\left[
\mathrm{softplus}(d_\theta(p)+m_{in})
\right].
\]

A good starting margin is \(m_{in}=0.01L_f\).

Then define pseudo-last support bands:

\[
\mathcal B_{\text{plantar}}
=
\{p\in S_L : h(p)\le \tau_p\},
\]

\[
\mathcal B_{\text{side}}
=
\{p\in S_L : \tau_p < h(p)\le \tau_s,\; |z-c(x)|>0.7a_t(x)\}.
\]

These are the hidden regions you most want to preserve: inner sole,
lower sidewall, heel cup.

Now softly anchor the shell near those bands:

\[
\mathcal L_{\text{band}}
=
\mathbb E_{p\sim \mathcal B_{\text{plantar}}\cup\mathcal B_{\text{side}}}
\left[
|d_\theta(p)+m_b|
\right],
\]

where \(m_b\) is a small desired shell offset. This prevents the
watertight shell from collapsing inward in the unseen sole/interior
zone.

\subsubsection{Then use the pseudo-last to classify which shell areas
should stay active in
mSDF}\label{then-use-the-pseudo-last-to-classify-which-shell-areas-should-stay-active-in-msdf}

G-Shell's key idea is to represent an open surface as an island on a
watertight shell via an mSDF. So after the shell exists, use the
pseudo-last to decide which shell neighborhoods are
\textbf{topologically valuable} and should be kept active.

Let \(q_i\) be shell vertices and \(m_\theta(q_i)\) the learned manifold
SDF value on the shell. Define a soft keep probability

\[
\hat t_i = \sigma\!\left(-\frac{m_\theta(q_i)}{\beta}\right),
\]

where lower mSDF means ``keep this shell area active.''

Now create pseudo-labels from visible shoe data and pseudo-last support
bands:

\[
t_i
=
\mathbf 1\Big[
d(q_i,S_{\text{vis}})<\tau_{\text{vis}}
\;\;\lor\;\;
d(q_i,\mathcal B_{\text{plantar}}\cup\mathcal B_{\text{side}})<\tau_L
\Big].
\]

Then train with

\[
\mathcal L_{\text{keep}}
=
-\frac{1}{N}\sum_i
\left[
t_i \log \hat t_i + (1-t_i)\log(1-\hat t_i)
\right].
\]

Add a shell-graph smoothness term to avoid fragmented islands:

\[
\mathcal L_{\text{smooth}}
=
\frac{1}{|E|}\sum_{(i,j)\in E}
w_{ij}\,|m_\theta(q_i)-m_\theta(q_j)|.
\]

The complete regularizer is then

\[
\mathcal L_{\text{pseudo-last}}
=
\lambda_c \mathcal L_{\text{contain}}
+\lambda_b \mathcal L_{\text{band}}
+\lambda_k \mathcal L_{\text{keep}}
+\lambda_s \mathcal L_{\text{smooth}}.
\]

In practice, I would use a curriculum:

\begin{itemize}
\tightlist
\item
  \textbf{early:} larger \(\lambda_c,\lambda_b\) so the shell forms
  around the last and visible outside;
\item
  \textbf{mid:} increase \(\lambda_k\) so the mSDF preserves
  support-related islands;
\item
  \textbf{late:} decay all pseudo-last weights and let image evidence
  dominate where it exists.
\end{itemize}

This is directly analogous in spirit to clothed-human reconstruction
pipelines that use a parametric body model to regularize unseen regions
without forcing the final surface to equal the template. PaMIR
conditions an implicit field on a parametric body model; ICON uses an
estimated SMPL body to guide occluded normal prediction and implicit
reconstruction; SeSDF explicitly deforms an SMPL-X-derived SDF to
recover the clothed surface.

\subsection{Failure cases and
validation}\label{failure-cases-and-validation}

The most common failure mode will be \textbf{bad foot-in-shoe
alignment}. If the aligned SUPR foot is too far forward, your toe
allowance becomes meaningless; if it is too low or too wide, the
pseudo-last will overconstrain the shell and hallucinate impossible
interior volume. This should be checked before any lofting or SDF
conversion.

The next failure mode is \textbf{using support width as if it were
actual inner width}. Many shoe meshes expose outsole/support shape more
reliably than the interior cavity, so directly using the raw support
width can overestimate the true inner width. That is why I recommend the
\(\eta_s\) clamp and why the pseudo-last should be a \textbf{soft
prior}, not a hard cavity template.

There are also failure modes from over-sculpting:

\begin{itemize}
\tightlist
\item
  too much toe merge makes the forefoot clownish;
\item
  too much arch preservation makes flat sneakers unrealistically hollow;
\item
  too much heel shrink creates self-intersections or an implausibly
  narrow backpart;
\item
  too much top inflation makes the shell over-thick and weakens
  visible-view supervision;
\item
  insufficient profile smoothing causes slice-to-slice ripples and
  unstable SDF signs.
\end{itemize}

Population variation is a quieter but real failure source. Large scan
studies and methodological reviews show foot-shape variation across sex
and regions, so one single fixed setting will not behave equally well
across all shoes and target feet. That is another reason to keep all
allowances normalized and learnable/tunable.

For validation, I would use a small but strict battery.

\textbf{Geometry diagnostics}

\begin{itemize}
\tightlist
\item
  \textbf{Plantar RMSE to footbed} \[
  \mathrm{RMSE}_{plantar}
  =
  \sqrt{
  \frac{1}{|S_b|}
  \sum_{(x,z)\in S_b}
  \left(y_b(x,z)-B(x,z)\right)^2
  }.
  \]
\item
  \textbf{Arch lift profile} along the centerline: \[
  a_{\text{arch}}(x)=B(x,c(x))-y_b(x,c(x)).
  \]
\item
  \textbf{Support conformity ratio} \[
  r_w(x)=\frac{2a_t(x)}{w(x)}.
  \] This should stay below your chosen \(\eta_s\) cap.
\item
  \textbf{Profile roughness} \[
  R_a=\sum_i (\Delta^2 a_i)^2,\qquad
  R_H=\sum_i (\Delta^2 H_i)^2.
  \] Sudden spikes indicate loft instability.
\end{itemize}

\textbf{Topology diagnostics}

\begin{itemize}
\tightlist
\item
  the mesh must be watertight,
\item
  no self-intersections,
\item
  exactly one connected toe-box component in slices after
  \(s_{\text{box}}\),
\item
  stable slice ordering and no contour flips.
\end{itemize}

\textbf{Containment diagnostics}

\begin{itemize}
\tightlist
\item
  fraction of pseudo-last surface points lying inside the coarse outer
  shoe envelope,
\item
  histogram of shell-to-last distances in the plantar and heel bands,
\item
  color map of clearance over the pseudo-last surface.
\end{itemize}

\textbf{Visualization diagnostics}

The single most useful visualization is a stack of sectional overlays
every 5\% of foot length. Plot, in each slice:

\begin{itemize}
\tightlist
\item
  raw SUPR section,
\item
  toe-merged/smoothed section,
\item
  pseudo-last section,
\item
  support interval \([z_L,z_R]\),
\item
  local footbed \(B\).
\end{itemize}

If those overlays look right, the 3D mesh almost always looks right. If
they do not, no amount of mesh smoothing will save you.

A broader caveat from the footwear-fit literature is that there is still
no universally accepted single metric of shoe fit, so section
diagnostics plus explicit dimension/error plots are more trustworthy
than a single scalar ``fit score.''

\subsection{Prior art and reading
list}\label{prior-art-and-reading-list}

The most relevant \textbf{shoe-last and foot-to-last} papers for your
project are \emph{A Deformation Method for Shoe Last Customization} by
Leng and Du, which uses aligned foot/last distance maps; \emph{A CAD
Approach for Designing Customized Shoe Last} by the same authors, which
explicitly adjusts foot data to toe spring and heel height before last
formation; \emph{Shoe-last Design Innovation for Better Shoe Fitting} by
Luximon and Luximon; \emph{3D Foot Scan to Custom Shoe Last} by Rout and
colleagues, which sections the foot at equal percentage intervals;
\emph{Computer Aided Design and Development of Customized Shoe Last} by
Sambhav and colleagues; \emph{Customization of shoe last based on 3D
design process with adjustable 3D ease allowance} by Mishra and
colleagues; and \emph{A Feature-based Shoe-last Design and its Automatic
Model Construction} by Li and colleagues. Those papers are the clearest
direct precedents for a section-based, parameterized pseudo-last
generator.

For \textbf{anatomical foot priors and foot-shape variability}, the most
relevant sources are \emph{SUPR: A Sparse Unified Part-Based Human
Representation}; \emph{Dynamic foot morphology explained through 4D
scanning and shape modeling}; \emph{Analysis of 1.2 million foot scans
from North America, Europe and Asia}; and the 2024 review \emph{Foot
shape assessment techniques for orthotic and footwear applications}.
Those provide the strongest justification for using SUPR as an
anatomical base while still smoothing it into a last-like envelope.

For \textbf{template/SDF priors in reconstruction}, the core references
are \emph{DeepSDF}, \emph{Occupancy Networks}, \emph{NeuS}, and
\emph{Ghost on the Shell}. Together they give you the right conceptual
ladder: continuous signed fields, continuous occupancy, multiview neural
surface reconstruction, and finally watertight-shell-plus-mSDF topology
control for open/non-watertight geometry.

For the \textbf{SMPL/clothed-human analogy} that is closest to your
proposed use of a pseudo-last, I would read \emph{PIFu}, \emph{PaMIR},
\emph{ICON}, and \emph{SeSDF}. The common theme is highly relevant to
your shoe problem: use an anatomical/template prior to stabilize
geometry in hidden or weakly observed regions, but learn a flexible
implicit field that is free to deviate where the observed surface
demands it.

The practical bottom line is that your pseudo-last should be implemented
as a \textbf{section-lofted, footbed-aware, toe-merged, heel-controlled,
watertight support prior}. If you do that, it becomes a very natural
supervisory object for both shell SDF containment and G-Shell mSDF
keep/cut decisions in the exact regions where turntable images give you
the least information.

\section{References and source links}\label{references-and-source-links}

\begin{enumerate}
\def\labelenumi{\arabic{enumi}.}
\tightlist
\item
  NeuS: Learning Neural Implicit Surfaces by Volume Rendering for
  Multi-view Reconstruction. \url{https://arxiv.org/abs/2106.10689}
\item
  G-Shell project page. \url{https://gshell3d.github.io/}
\item
  Ghost on the Shell / G-Shell paper page.
  \url{https://openreview.net/forum?id=Ad87VjRqUw}
\item
  SUPR: A Sparse Unified Part-Based Human Representation.
  \url{https://arxiv.org/abs/2210.13861}
\item
  A Deformation Method for Shoe Last Customization.
  \url{https://www.cad-journal.net/files/vol_2/CAD_2\%281-4\%29_2005_11-18.pdf}
\item
  A CAD Approach for Designing Customized Shoe Last.
  \url{https://www.cad-journal.net/files/vol_3/CAD_3\%281-4\%29_2006_377-384.pdf}
\item
  Shoe-last Design Innovation for Better Shoe Fitting.
  \url{https://www.sciencedirect.com/science/article/abs/pii/S0166361509001365}
\item
  3D Foot Scan to Custom Shoe Last.
  \url{https://www.researchgate.net/publication/228879419_3D_Foot_Scan_to_Custom_Shoe_Last}
\item
  Computer Aided Design and Development of Customized Shoe Last.
  \url{https://www.cad-journal.net/files/vol_8/CAD_8\%286\%29_2011_819-826.pdf}
\item
  Customization of shoe last based on 3D design process with adjustable
  3D ease allowance for better comfort and design.
  \url{https://link.springer.com/article/10.1007/s00170-022-10427-5}
\item
  A Feature-based Shoe-last Design and its Automatic Model Construction.
  \url{https://cad-journal.net/files/vol_22/CAD_22\%284\%29_2025_600-615.pdf}
\item
  Dynamic foot morphology explained through 4D scanning and shape
  modeling. \url{https://arxiv.org/abs/2007.11077}
\item
  Analysis of 1.2 million foot scans from North America, Europe and
  Asia. \url{https://www.nature.com/articles/s41598-019-55432-z}
\item
  Foot shape assessment techniques for orthotic and footwear
  applications.
  \url{https://www.frontiersin.org/journals/bioengineering-and-biotechnology/articles/10.3389/fbioe.2024.1416499/full}
\item
  Robust estimation of foot dimensions from 3D foot scan data.
  \url{https://proc.3dbody.tech/papers/A2012/a12147_20son.pdf}
\item
  Quantitative comparison of plantar foot shapes under different
  weight-bearing conditions.
  \url{https://www.researchgate.net/publication/8624319_Quantitative_comparison_of_plantar_foot_shapes_under_different_weight-bearing_conditions}
\item
  Plantar foot shapes under different weight-bearing conditions.
  \url{https://ira.lib.polyu.edu.hk/bitstream/10397/6811/1/Tsung_Plantar_Foot_Shapes.pdf}
\item
  The effect of shoe toe box shape and volume on forefoot interdigital
  and plantar pressures in healthy females.
  \url{https://pubmed.ncbi.nlm.nih.gov/23886242/}
\item
  Research profile page for toe-box pressure paper.
  \url{https://researchprofiles.herts.ac.uk/en/publications/the-effect-of-shoe-toe-box-shape-and-volume-on-forefoot-interdigi/}
\item
  Does the shoe really fit? Characterising ill-fitting footwear among
  community-dwelling older adults attending podiatry services.
  \url{https://pmc.ncbi.nlm.nih.gov/articles/PMC7020372/}
\item
  Comparison of shoe-length fit between people with and without diabetic
  peripheral neuropathy.
  \url{https://pmc.ncbi.nlm.nih.gov/articles/PMC3407779/}
\item
  An exploratory study of dynamic foot shape measurements.
  \url{https://pmc.ncbi.nlm.nih.gov/articles/PMC10224935/}
\item
  A methodology for determining the allowances for fitting footwear.
  \url{https://www.researchgate.net/publication/264438902_A_methodology_for_determining_the_allowances_for_fitting_footwear}
\item
  Foot deformation during walking.
  \url{https://www.researchgate.net/publication/261254608_Foot_deformation_during_walking_Differences_between_static_and_dynamic_3D_foot_morphology_in_developing_feet}
\item
  Dynamic foot shape and safety footwear study.
  \url{https://www.tandfonline.com/doi/full/10.1080/19424280.2018.1529062}
\item
  Generalized winding numbers.
  \url{https://dl.acm.org/doi/10.1145/2461912.2461916}
\item
  Fast winding numbers for soups and clouds.
  \url{https://www.dgp.toronto.edu/projects/fast-winding-numbers/}
\item
  Fast winding numbers paper PDF.
  \url{https://www.dgp.toronto.edu/projects/fast-winding-numbers/fast-winding-numbers-for-soups-and-clouds-siggraph-2018-compressed-barill-et-al.pdf}
\item
  DeepSDF: Learning Continuous Signed Distance Functions for Shape
  Representation. \url{https://arxiv.org/abs/1901.05103}
\item
  Occupancy Networks: Learning 3D Reconstruction in Function Space.
  \url{https://arxiv.org/abs/1812.03828}
\item
  PIFu: Pixel-Aligned Implicit Function for High-Resolution Clothed
  Human Digitization. \url{https://arxiv.org/abs/1905.05172}
\item
  PaMIR: Parametric Model-Conditioned Implicit Representation for
  Image-based Human Reconstruction.
  \url{https://arxiv.org/abs/2007.03858}
\item
  ICON project page. \url{https://icon.is.tue.mpg.de/}
\item
  ICON: Implicit Clothed Humans Obtained from Normals.
  \url{https://openaccess.thecvf.com/content/CVPR2022/papers/Xiu_ICON_Implicit_Clothed_Humans_Obtained_From_Normals_CVPR_2022_paper.pdf}
\item
  SeSDF: Self-Evolved Signed Distance Field for Implicit 3D Clothed
  Human Reconstruction.
  \url{https://openaccess.thecvf.com/content/CVPR2023/papers/Cao_SeSDF_Self-Evolved_Signed_Distance_Field_for_Implicit_3D_Clothed_Human_CVPR_2023_paper.pdf}
\item
  Footwear comfort: a systematic search and narrative synthesis.
  \url{https://pmc.ncbi.nlm.nih.gov/articles/PMC8650278/}
\item
  Footwear Toe-Box Shape and Medial Forefoot Pressures.
  \url{https://pmc.ncbi.nlm.nih.gov/articles/PMC11992996/}
\item
  Medical-grade footwear: the impact of fit and comfort.
  \url{https://pmc.ncbi.nlm.nih.gov/articles/PMC5217416/}
\item
  Toe Box Shape of Running Shoes Affects In-Shoe Foot Function.
  \url{https://www.mdpi.com/2306-5354/11/5/457}
\item
  Human foot modeling towards footwear design.
  \url{https://www.sciencedirect.com/science/article/abs/pii/S001044851100193X}
\item
  NeuS2 project page. \url{https://vcai.mpi-inf.mpg.de/projects/NeuS2/}
\item
  Learning Signed Distance Field for Multi-view Surface Reconstruction.
  \url{https://arxiv.org/abs/2108.09964}
\item
  Distance Field Rasterization for End-to-End Mesh Reconstruction.
  \url{https://arxiv.org/html/2604.23537v1}
\item
  GPU Gems Chapter 34: Signed Distance Fields Using Single-Pass GPU Scan
  Conversion of Tetrahedra.
  \url{https://developer.nvidia.com/gpugems/gpugems3/part-v-physics-simulation/chapter-34-signed-distance-fields-using-single-pass-gpu}
\item
  Automated Shoe Last Customization using MATLAB.
  \url{https://www.scitepress.org/Papers/2019/79493/79493.pdf}
\item
  Laser Line-Scanning Based Customized Shoe-Last System.
  \url{https://www.researchgate.net/publication/251970370_Laser_Line-Scanning_Based_Customized_Shoe-Last_System}
\item
  Rapid parametric design methods for shoe-last customization.
  \url{https://www.researchgate.net/publication/227333110_Rapid_parametric_design_methods_for_shoe-last_customization}
\item
  Morphing-Based Surface Blending Operator for Footwear CAD.
  \url{https://www.researchgate.net/publication/267586204_A_Morphing-Based_Surface_Blending_Operator_for_Footwear_CAD}
\item
  Foot shape prediction using elliptical Fourier analysis.
  \url{https://www.researchgate.net/publication/313827280_Foot_shape_prediction_using_elliptical_Fourier_analysis}
\item
  3D foot shape generation from 2D information.
  \url{https://www.researchgate.net/publication/7671857_3D_foot_shape_generation_from_2D_information}
\end{enumerate}

\end{document}
